% Source-aligned English for printed p.151.
\SourcePage{151}

\LayerTarget{tohoku:construction:resolving-functor-spectral-sequences:p2}
\noindent hyperhomology spectral sequences (2.4.2) of $F$ with respect to the
complex $K$.

\LayerTarget{tohoku:prop:2.5.3}
\noindent\textsc{Proposition 2.5.3.} --- \textit{The isomorphisms (2.5.3)
arise from natural isomorphisms of the first, respectively second, spectral
sequence of the bicomplex $\mathbf F(K)$ with the first, respectively second,
hyperhomology spectral sequence of $F$ with respect to $K$.}

\LayerTarget{tohoku:prop:2.5.3:proof:p1}
We now make these isomorphisms explicit. Let
$C(K)=(C^{p,q}(K))_{p,q}$ be an \textquotedblleft injective
resolution\textquotedblright{} of $K$ in the sense of [6, Chapter XVII].
Regard $K$ as a bicomplex whose second degree is zero; there is then an
augmentation homomorphism $K\longrightarrow C(K)$. Consider the bicomplex
$M(K)=\mathbf F(C(K))$ defined by
\[
M^{p,q}(K)=\sum_{q'+q''=q}\mathbf F^{q''}(C^{p,q'}(K)).
\]
Although $M(K)$ is not uniquely determined by $K$, its two spectral sequences
and their common abutment are: $C(K)$ is determined up to a homotopy
equivalence of bicomplexes. They are therefore well-defined spectral functors
of $K$. Put $L(K)=F(C(K))$. The augmentations
$K\longrightarrow C(K)$ and $F\longrightarrow\mathbf F$ define
homomorphisms of bicomplexes
\[
L(K)=F(C(K))\longrightarrow M(K)=\mathbf F(C(K))
\longleftarrow\mathbf F(K).
\]
Thus there are natural homomorphisms of the corresponding spectral sequences,
independent of the choice of $C(K)$:
\LayerTarget{tohoku:eq:2.5.4}
\[
\tag{2.5.4}
\begin{aligned}
\mathrm I L(K)&\longrightarrow\mathrm I M(K)\longleftarrow\mathrm I\mathbf F(K),\\
\mathrm{II}L(K)&\longrightarrow\mathrm{II}M(K)\longleftarrow\mathrm{II}\mathbf F(K).
\end{aligned}
\]
They define the same homomorphisms on the abutments,
\[
HL(K)\longrightarrow HM(K)\longleftarrow H\mathbf F(K).
\]
On the initial terms, the homomorphisms in (2.5.4) give
\LayerTarget{tohoku:eq:2.5.5}
\[
\tag{2.5.5}
\begin{aligned}
{}_{\mathrm I}E_2^{p,q}L(K)&\longrightarrow{}_{\mathrm I}E_2^{p,q}M(K)
\longleftarrow{}_{\mathrm I}E_2^{p,q}\mathbf F(K),\\
{}_{\mathrm{II}}E_2^{p,q}L(K)&\longrightarrow{}_{\mathrm{II}}E_2^{p,q}M(K)
\longleftarrow{}_{\mathrm{II}}E_2^{p,q}\mathbf F(K).
\end{aligned}
\]
We shall prove that the homomorphisms in (2.5.5) are isomorphisms, and that
the resulting isomorphisms between the extreme terms of (2.5.5), which are
the initial terms of the hyperhomology spectral sequences of $F$ for $K$ and
of the bicomplex $\mathbf F(K)$, are those in (2.5.3). It will follow that
the homomorphisms in (2.5.4) are isomorphisms, and the resulting isomorphisms
between their extreme terms are the ones sought. It therefore suffices to
show that the middle terms in the two lines of (2.5.5) have the forms
$H^p(R^q(F(K)))$ and $R^pF(H^q(K))$, respectively. Once this is done,
checking that the homomorphisms in (2.5.5) are the natural isomorphisms is
purely mechanical and is implicit in the proof that follows.

We shall show that $H_{\mathrm{II}}(M)^{p,q}=R^qF(K^p)$, which immediately
gives $ {}_{\mathrm I}E_2^{p,q}(M)=H^p(R^qF(K))$. For fixed $p$,
$M^{p*}=(M^{p,q})_q$ is the simple complex associated to the bicomplex
\[
(\mathbf F^{q''}(C^{p,q'}(K)))_{q',q''}=\mathbf F(C(K^p)),
\]
where $C(K^p)$ denotes the complex $C^{p*}(K)$,

\TohokuContinuation{152}

